\newpage
\phantomsection
\label{prf:restriction-preserves-monotonicity}
\LRAProofFor{prop:restriction-preserves-monotonicity}

\begin{remark*}[Return]
\hyperref[prop:restriction-preserves-monotonicity]{Return to Proposition}
\end{remark*}

\begin{proposition*}[Restriction Preserves Monotonicity]
Restricting the domain of an order-preserving or order-reversing map preserves
that monotonicity type.
\end{proposition*}

\begin{proof}[Professional Standard Proof]
\LRAProofBodyStart
Let $c_1,c_2\in C$ with $c_1\leq_C c_2$. By the induced order,
$c_1\leq c_2$ in $P$. If $f$ is order-preserving, then
$f(c_1)\leq' f(c_2)$. If $f$ is order-reversing, then
$f(c_2)\leq' f(c_1)$.
\end{proof}

\begin{proof}[Detailed Learning Proof]
\LRAProofBodyStart
The induced order on $C$ uses the same comparisons inherited from $P$.
Therefore the same monotonicity implication that worked for $f$ on $P$ also
works for the restricted map on $C$.
\end{proof}
\begin{remark*}[Proof structure]
\end{remark*}



\begin{dependencies}
\begin{itemize}
  \item \hyperref[prop:restriction-preserves-monotonicity]{Proof target}
\end{itemize}
\end{dependencies}

\clearpage

